\begin{textsample}{POS Dim 2 – ma – Score 31.00 – ma/ma\_em\_25.txt}  \label{ex:f2_pos_046}
2.
A organização curricular do ensino médio para o estado do Maranhão Ao situar o ensino médio dentro da atual estrutura da educação brasileira, deve -se \textbf{reiterar} que ela advém de mudanças legais ocorridas ao longo da história da educação do país, regulamentadas pela Lei de \textbf{Diretrizes} e Bases da Educação ( LDB ) e suas respectivas alterações ; e esta, por sua vez, vem cumprir as determinações gerais da Constituição Federal de 1988.
O ensino médio, conforme o \textbf{art}. 21, é uma etapa específica da educação básica : de acordo com o \textbf{art}. 21 da LDB, a educação escolar compõe -se de : 1.
Educação básica, formada pela educação infantil, ensino fundamental e ensino médio ; 2.
Educação superior.
Ainda conforme a mesma \textbf{legislação}, a educação básica “ tem por finalidade desenvolver o educando, assegurar -lhe a formação comum indispensável para o exercício da cidadania e fornecer -lhe meios para progredir no trabalho e em estudos posteriores ” ( \textbf{art}. 22 ).
Sua \textbf{oferta} será garantida no ensino regular e suas diferentes possibilidades, contemplando todas as modalidades de ensino : educação de jovens e adultos, educação especial e educação profissional, sendo que esta última também é uma modalidade da educação superior.
Dessa forma, a \textbf{oferta} do ensino médio deve \textbf{atender} ao direito subjetivo de cada cidadão e ser assegurada a todos e todas, em cada escola do país e, do mesmo modo, em cada espaço do território maranhense.
Sendo assim, no compromisso com a \textbf{oferta} da educação pública, o estado do Maranhão atende toda sua diversidade de público : quilombolas, indígenas, campo, apenados, EJA e educação especial, tendo, ainda, sido garantida a expansão do acesso e permanência.
Também atende com a ampliação das escolas de tempo integral, por meio dos 55 Centros Educa Mais, que gradativamente estão sendo ampliadas, buscando -se alcançar a meta da \textbf{oferta} da escola em tempo integral até 2024 \textbf{prevista} no Plano Nacional de Educação e \textbf{reiterada} no Plano Estadual de Educação do Maranhão ( PEE -MA ).
Destaca -se a importância de reconhecer os desafios próprios do ensino médio que se apresentam nas redes públicas de ensino no país para que possam ser planejadas ações que venham, se não a eliminar, pelo menos a minimizar os entraves para a \textbf{oferta} de uma educação pública de qualidade.
Para o enfrentamento de tais desafios, o \textbf{governo} do estado do Maranhão, por meio da \textbf{política} educacional \textbf{instituída} com a Escola Digna, realizou um planejamento que contempla intervenções desde a estrutura física das escolas da rede pública estadual e \textbf{municipal}, com formação de professores, investimentos em tecnologia e outros recursos pedagógicos.
Hoje, é possível, para o estado, apresentar resultados positivos com foco em metas previamente definidas para toda a rede estadual.
Ressalta -se que, para todo esse planejamento, foi realizado o \textbf{diagnóstico} da rede estadual, que \textbf{culminou} na definição das \textbf{políticas} educacionais do atual \textbf{governo} e que teve origem nas escutas pedagógicas realizadas junto a gestores, docentes e discentes de todas as 19 Unidades Regionais de Educação, dando, assim, o direcionamento para todas as ações do macroplanejamento da Secretaria de Estado da Educação.
Outro instrumento indispensável para a \textbf{elevação} da qualidade da \textbf{oferta} do ensino público é o Sistema de Avaliação Estadual do Maranhão ( Seama ), que faz as avaliações anuais.
Estas avaliações possibilitam a identificação de fragilidades, ao detectar onde é necessária maior atenção junto aos discentes, apoiando o monitoramento das aprendizagens de forma mais pontual.
O estado do Maranhão ( Censo Escolar 2021 ) apresenta 335.629 estudantes matriculados no ensino médio, que são \textbf{atendidos} por 1.156 escolas.
Destas, 787 estão localizadas na zona urbana e 369 escolas, na zona rural.
Neste universo, temos, ainda, escolas que \textbf{ofertam} educação profissional e que funcionam em tempo integral.
Para que o desenvolvimento do ensino médio funcione, a rede apresenta 14.062 professores no ensino regular, 3.365 professores na educação profissional ( médio integrado, FIC e Técnico ) e 4.117 professores na educação de jovens e adultos ( EJA ).
Além disso, tem -se os dados dos itens 2.1 e 2.2, que trazem os desafios a serem superados para a \textbf{elevação} da aprendizagem escolar, demonstrando que muito se avançou, mas que ainda há muito a ser feito.
Esses dados indicam, claramente, onde é preciso investir mais para a superação dos desafios.
Todo o trabalho de monitoramento e acompanhamento da rede escolar é realizado por meio do Seama e do Siaep, das orientações curriculares e das \textbf{diretrizes} curriculares do estado do Maranhão para o ensino médio e, ainda, por meio de \textbf{cadernos} pedagógicos de estudantes e professores, inclusive com material de apoio escolar aos estudantes que ficam de pendência em sua aprovação para o ano seguinte de sua escolarização, que obedecem a \textbf{legislação} específica do Conselho Estadual de Educação do Maranhão.
Diante da obrigatoriedade de implantação do Novo Ensino Médio no ano de 2020 e da relevante iniciativa da implantação e expansão da educação integral, por meio dos Centros Educa Mais, foram iniciadas as ações referentes ao Novo Ensino Médio, com base na Lei no 13.415, de 16 de \textbf{fevereiro} de 2017, que aprovou a \textbf{reforma} do ensino médio e alterou a Lei de \textbf{Diretrizes} e Bases, bem como as \textbf{Diretrizes} Curriculares Nacionais para o Ensino Médio, \textbf{atualizadas} por meio da Resolução CNE/CEB no 3, de 21 de \textbf{novembro} de 2018, que determina, no \textbf{art}. 4o : “ As \textbf{instituições} de ensino que \textbf{ofertam} essa etapa da Educação Básica devem estruturar suas propostas pedagógicas considerando as finalidades \textbf{previstas} no \textbf{art}. 35 da Lei no 9.394/1996, de \textbf{Diretrizes} e Bases da Educação Nacional ”.
Logo, a partir dessas \textbf{deliberações}, deu -se a construção de uma nova proposta curricular pelo estado do Maranhão, que norteará as propostas pedagógicas de cada unidade escolar.
Vale destacar que a \textbf{legislação} citada traz a concretização da flexibilização curricular, que já havia sido mencionada na LDB, mas que, hoje, é possibilitada por meio de outros instrumentos legais que não só regulamentam, mas orientam todo o processo de flexibilização curricular trazida pela lei do ensino médio, sem deixar de ser considerada a autonomia das redes e das \textbf{instituições} de ensino, além das especificidades locais que as envolvem.
A experiência acumulada pela Secretaria de Educação com a implantação e o desenvolvimento da \textbf{política} de educação integral maranhense desde 2015 permitiu que os seus aprendizados fossem expandidos e aplicados na concepção do Novo Ensino Médio como \textbf{política} pública que assegura um único \textbf{currículo} para os Centros de Ensino de tempo parcial e Centros Educa Mais.

% matched lemmas: art, atender, atualizar, caderno, culminar, currículo, deliberação, diagnóstico, diretriz, elevação, fevereiro, governo, instituir, instituição, legislação, municipal, novembro, oferta, ofertar, política, prever, reforma, reiterar
\end{textsample}
